\section*{Wstęp}
Paweł Cicheńki z Kępy Bazarowej jest osobą o szerokich zainteresowaniach, łączącą sport, technikę i podróże. Poniższy opis przedstawia jego umiejętności oraz sposób, w jaki spędza wolny czas. Dokument może służyć jako krótka prezentacja kompetencji.

\section*{Umiejętności techniczne}
\begin{tabularx}{\textwidth}{@{}lX@{}}
\toprule
\textbf{Obszar} & \textbf{Opis} \\
\midrule
Elektryka & Dobre umiejętności elektryczne \\
Spawanie i robotyka & Wyspawał zaawansowanego robota minisumo \\
Informatyka & Student informatyki w Wielkopolsce \\
\bottomrule
\end{tabularx}

\vspace{0.4cm}

\section*{Opis kompetencji}
Paweł jest dobrym elektrykiem, co potwierdza praktycznym projektem, jakim był zaawansowany robot minisumo. Konstrukcja ta wymagała od niego zarówno umiejętności spawania, jak i znajomości zagadnień elektrycznych. Dzięki temu łączy wiedzę praktyczną z kierunkiem studiów, którym jest informatyka.

Wyjazd na studia do Wielkopolski pozwolił mu rozwijać się w nowym środowisku, bez rezygnacji z pasji sportowych. Nadal kocha lekkoatletykę i pozostaje z nią związany.

\section*{Zainteresowania}
\begin{itemize}
\item Lekkoatletyka, którą kocha
\item Podróże po świecie, ostatnio do Kenii i Berlina
\item Wędkarstwo, w tym polowanie na ryby
\item Kajakarstwo
\end{itemize}

\section*{Czas wolny}
W wolnym czasie Paweł poluje na ryby i bierze udział w lokalnych konkursach wędkarskich, w jednym z nich zajął 12. miejsce. Bardzo lubi też kajakować, a jego mieszkanie położone jest nad jeziorem, co sprzyja obu tym zajęciom.

\section*{Podróże i kontakty}
Paweł podróżuje po świecie. Ostatnio odwiedził Kenię oraz Berlin, gdzie spotkał się z licznymi CEO. Doświadczenia te poszerzają jego perspektywę i pozwalają poznawać różne środowiska.

\section*{Podsumowanie}
Paweł Cicheńki to osoba łącząca sport, technikę i podróże. Jego umiejętności elektryczne i konstruktorskie uzupełniają studia informatyczne, a pasje sportowe i wodne dopełniają obraz aktywnego stylu życia.